\section{Introduction}
The interaction between the judiciary and policymaking is controversial. For instance, in the regulatory context, courts and agencies might be seen as opposing agents in terms of markets’ social control \citep{glaeser2003rise}. Agencies may be more efficient at resolving conflicts than courts, which are often viewed as more expensive, unpredictable, and biased.

The execution of contracts in the context of public procurement has led to some results related to conflict resolution. The judiciary may, directly or indirectly, influence public tender outcomes. For instance, inefficient courts may induce public buyers to apply penalties instead of solving the conflict through litigation, leading to longer delays in public works deliveries and payments in litigated contracts \citep{coviello2017court}.

Health litigation is another example of the judiciary’s interference in the executive branch \citep{wang2015right}. This paper aims to estimate the enforcement costs of health litigation and administrative requests. In other words, the main objective is to assess the government waste generated when the judiciary directly affects public policy.

The Brazilian Constitution states that all citizens have the “right to life,” and the state has the explicit public health objective of “providing universal coverage” for everyone. Available or potential budget resources must be considered across time for these objectives to be materialized and sustainable in the long run.

The Unified Health System (SUS) is precisely the materialization of this statement: a coordinated set of financially viable actions, public programs, and infrastructure, aiming to achieve the objectives established in the Constitution. Based on aggregate social preferences, the government chooses priorities and implements public policies subject to budget restrictions.

However, the judiciary has a strict interpretation of the Constitution that generally ignores the budgetary dimension. This means that in Brazil, it is possible to obtain any medication or medical procedure through litigation regardless of the costs involved. Court orders have granted a range, for example, from acetylsalicylic acid (aspirin or similar) to galsulfase, indicated for treating rare and severe joint disease (mucopolysaccharidosis type VI). Individual treatment with galsulfase has an estimated annual cost of US\$400 thousand.

This rigid interpretation of law leads to significant distortions in implementing health policies, such as public procurement of prescription and nonprescription drugs. Court decisions are enforced as preliminary injunctions that require the government to make purchases within one-third of the time planned acquisitions are made, hampering all public buying processes.

The way the planning procedures for the acquisition of these goods are carried out may substantially affect the procurement conditions or outcomes and ultimately might influence the results of public policies. Public bureaus should have the appropriate time to identify all needs and ends, analyze market conditions, and set relevant tender parameters (item specifications, quantities, and reference prices, for instance).

Favorable planning circumstances may increase the likelihood of achieving public policy goals efficiently and effectively. On the other hand, under unfavorable planning conditions, purchases might be inadequate to meet public needs and more expensive, undermining public policies’ final impact.

There is a vast literature on the waste of resources in public services, including public procurement. A prominent approach separates the causes of waste into two primary sources: corruption and mismanagement \citep{bandiera2009active}.

The involvement of officials in corruption (i.e., active waste), such as favoring private firms in public tenders in exchange for bribery, has received much attention in the literature and from policymakers for its impact on public procurement efficiency \citep{mironov2016corruption, basheka2011economic}.

On the other hand, mismanagement (passive waste) might lead to higher prices for various other reasons, such as inadequate civil servant skills \citep{best2017individuals}, lack of incentives to minimize costs \citep{cullen2016performance, ashraf2016dogooders, bandiera2017motivating} or improper management practices \citep{williams2018management, bloom2015impact, kvasnicka2015auctions, lewisfaupel2016can, rasul2018management}, which may increase the probability of collusion or bid rigging \citep{clark2018bid, moore2012cartels}.

Mismanagement and corruption are often associated mainly with internal aspects of public administration. However, the external dimension is quite relevant to understanding the functioning and distortions of the procurement process. Restrictions imposed and behaviors shown by control agencies and other stakeholders acting as watchdogs might strongly influence public officials’ decisions and, as a result, undermine efficiency \citep{brewer2010impact}.

Accordingly, this paper contributes to the vast literature on public procurement efficiency \citep{bandiera2009active, ashraf2016dogooders, lewisfaupel2016can} by providing evidence that external shocks may affect procurement outcomes by harming the \textit{ex ante} process of planning. Notably, this policy experiment allows the inefficiency due to the judicial review of public procurement to be estimated.

Health litigation and administrative requests are exogenous shocks that affect how the government buys prescription and over-the-counter drugs. Thus, those shocks can be separated into urgent (litigated and administrative requests) and ordinary (standard procedure) types of purchases, treatments, and control groups, respectively.

We estimate the impact of health litigation and administrative requests (planning and executing a tender) on public procurement efficiency, comparing urgent and ordinary purchases. The objective is to assess the effects of those court orders’ enforcement costs and administrative demands on public tender efficiency.

Additionally, we compare litigated and administrative purchases to identify the “under the gun” effect. Administrative and litigated purchases are similar in all adverse planning and execution conditions. However, if the government fails to comply with a court order to purchase medicines, public bureaus are subject to severe punishment. Thus, the “under the gun” effect is an attempt to isolate the possibility of severe penalties as an additional cost to the government.

We construct unique administrative data on bid-level public procurement transactions of litigated, administrative, and ordinary health-related item purchases in the state of São Paulo, Brazil, from January 2009 to December 2019.

The main findings indicate higher reference prices for urgent than ordinary purchases (from 60.4\% to 68.93\% higher). This result suggests that unfavorable conditions for compliance with court orders or administrative requests (shorter delivery time, lower quantities, and the threat of punishment) significantly increase expectations regarding acquisition costs.

Moreover, the (over)enforcement costs consist of (i) higher negotiated prices (from 30.73\% to 44.37\% higher), (ii) fewer participant firms (from 28.63\% to 32.21\% fewer), (iii) fewer bids (from 39.40\% to 45.93\% fewer), and (iv) a lower probability of success (from 38.56\% to 48.66\% less probable) in urgent tenders than in ordinary purchases.

Finally, we estimate the “under the gun” effect: a litigated purchase is between 8.83\% and 9.97\% more expensive than an administrative request, a difference attributed to the possibility of a judicial punishment of government members in the first and not in the second case.

In summary, judicial decisions compel the executive branch to carry out purchases that generate high public budget costs.